\documentclass[11pt, a4paper]{problems_template}

\begin{document}

\begin{titlepage}
  \begin{center}
	\Huge{Санкт-Петербургская математическая олимпиада}
  \end{center}
\end{titlepage} 

\chapter{1961}

\begin{exercise}[21]
Докажите, что шахматную доску $10 \times 10$ нельзя покрыть фигурками $T$-тетрамино.
\end{exercise}

\begin{Solution}
Рассмотрим шахматную раскраску доски. Каждая фигурка покрывает либо 3 черных клетки и одну белую либо 3 белых клетки и одну черную.  Обозначим через $a$ число фигур первого типа, а через $b$ - второго. Тогда:
$$
3a + b = 50, \quad 3b + a = 50, \quad a + b = 25
$$
Откуда получается, что $a = b$ и $a + b = 25$. Противоречие.
\end{Solution}

\begin{exercise}[23]
Докажите, что все числа вида $1156, 111556, 11115556, \dots$ являются точными квадратами.
\end{exercise}

\begin{Solution}
Легко заметить, что все числа этого вида представляются в форме
$$
\frac{10^{n}-1}{9} \cdot 10^{n} + \frac{10^{n} - 1}{9} \cdot 5 + 1 = \frac{10^{2n} + 4 \cdot 10^{n} + 4}{9} = \frac{(10^{n} + 2)^{2}}{3^{2}}
$$
\end{Solution}

\chapter{1962}

\begin{exercise}[5]
Докажите, что шахматную доску $201 \times 201$ можно обойти ходом шахматного коня, побывав на каждом поле ровно один раз.
\end{exercise}

\begin{exercise}[32]
Освободитесь от иррациональности в знаменателе
$$
\frac{1}{\sqrt[3]{a} + \sqrt[3]{b} + \sqrt[3]{c}}
$$
\end{exercise}

\begin{Solution}
Начнем с тождества
$$
(x^{2} + y^{2} + z^{2} - xy - yz - xz)(x + y + z) = x^{3} + y^{3} + z^{3} - 3 xyz
$$
Тогда
$$
\frac{1}{x+y+z} = \frac{x^{2} + y^{2} + z^{2} - xy - yz - xz}{x^{3} + y^{3} + z^{3} - 3xyz} = \frac{(x^{2} + y^{2} + z^{2} - xy - yz - xz)((x^{3} + y^{3} + z^{3})^{2} + (3xyz)^{2} + 3xyz(x^{3} + y^{3} + z^{3}))}{(x^{3} + y^{3} + z^{3})^{3} - (3xyz)^{3}}
$$
Подставляя $x = \sqrt[3]{a}$, $y = \sqrt[3]{b}$, $z = \sqrt[3]{c}$
$$
\frac{1}{\sqrt[3]{a} + \sqrt[3]{b} + \sqrt[3]{c}} = \frac{(\sqrt[3]{a^{2}} + \sqrt[3]{b^{2}} + \sqrt[3]{c^{2}} - \sqrt[3]{ab} - \sqrt[3]{ac} - \sqrt[3]{bc})((a + b + c)^{2} + 9\sqrt[3]{a^{2}b^{2}c^{2}} + 3\sqrt[3]{abc}(a+b+c))}{(a+b+c)^{3} - 27abc}
$$
\end{Solution}

\begin{exercise}[37]
Дан многогранник. Число сторон всех его граней, кроме одной, делится на данное натуральное $n$, большее 1. Докажите, что грани этого многогранника нельзя раскрасить двумя красками так, чтобы соседние грани были окрашены в разные цвета.
\end{exercise}

\begin{Solution}
Пусть многогранник окрашен в черный и белый цвет. Пусть единственная грань, число сторон которой не делится на $n$, будет окрашена в белый. Посчитаем количество ребер соединяющих черную и белую грани. С одной стороны
$$
E = \sum_{i = 1}^{N_{w}} n\cdot q_{i} + q = q \mod n
$$
С другой стороны
$$
E = \sum_{i = 1}^{N_{b}} n\cdot p_{i} = 0 \mod n
$$
Противоречие.
\end{Solution}

\begin{exercise}[54]
Рассматриваются бесконечные последовательности, состоящие из чисел $1, 2, \dots, N$. Два члена в двух \\ последовательностях назовем одинаковыми, если они равны и имеют одинаковые номера. Сколько можно составить последовательностей так, чтобы любые две из них имели ровно $k$ одинаковых членов?
\end{exercise}


\chapter{1965}

\begin{exercise}[26]
Найдите сумму
$$
\frac{1}{\cos \alpha \cos 2 \alpha} + \frac{1}{\cos 2\alpha \cos 3 \alpha} + \dots + \frac{1}{\cos (n - 1)\alpha \cos n\alpha}
$$
\end{exercise}

\begin{Solution}
Рассмотрим разность
$$
\tg (n+1)\alpha - \tg n\alpha = \frac{\sin(n+1)\alpha}{\cos(n+1)\alpha} - \frac{\sin n\alpha}{\cos n\alpha} = \frac{\sin(n+1)\alpha \cdot \cos n\alpha -\cos(n+1)\alpha \cdot \sin n\alpha}{\cos(n+1)\alpha\cos n\alpha} = \frac{\sin \alpha}{\cos(n+1)\alpha\cos n\alpha}
$$
Тогда исходная сумма представляется в виде
$$
\frac{\tg 2\alpha - \tg \alpha + \tg 3\alpha - \tg 2\alpha + \dots + \tg n\alpha - \tg (n-1)\alpha}{\sin \alpha} = \frac{\tg n\alpha - \tg \alpha}{\sin\alpha} = \frac{\sin(n-1)\alpha}{\cos\alpha\cos n\alpha\sin\alpha} = \frac{2\sin(n-1)\alpha}{\sin 2\alpha \cos n\alpha}
$$
\end{Solution}

\chapter{1967}

\begin{exercise}
В компании 18 человек. Докажите, что среди них есть четверо попарно незнакомых людей или среди них есть четверо попарно знакомых.
\end{exercise}

\chapter{1968}

\begin{exercise}[36]
Найдите такое целое $n$, что среди цифр десятичной записи числа $5^{n}$ есть по крайней мере 1968 нулей подряд.
\end{exercise}

\chapter{1969}

\begin{exercise}[32]
На площади собралось 50 гангстеров. Они одновременно стреляют друг в друга, причем каждый стреляет в ближайшего к нему или в одного из ближайших. Каково минимальное возможное количество убитых?
\end{exercise}

\chapter{1970}

\begin{exercise}[16]
В квадрате $5 \times 5$ закрашено 16 клеток. Докажите, что в нем можно выбрать квадрат $2 \times 2$, в котором закрашено не менее трех клеток.
\end{exercise}

\begin{exercise}[22]
Положительные числа $a$, $b$, $c$, $d$ таковы, что $a^{2} + b^{2} = c^{2} + d^{2}$ и $a^{3} + b^{3} = c^{3} + d^{3}$. Докажите, что $ab = cd$.
\end{exercise}

\begin{Solution}

\end{Solution}


\begin{exercise}[34]
Натуральные числа $a_{1}, a_{2}, \dots$, таковы, что $0 < a_{1} < a_{2} < \dots$ и $a_{n} < 2n$ для любого $n$. Докажите, что любое натуральное число можно представить в виде разности двух чисел из этой последовательности или как число из самой последовательности.
\end{exercise}

\chapter{1971}

\begin{exercise}[14]
Даны числа $5^{1971}$ и $2^{1971}$. Они записаны подряд. Каково количество цифр полученного числа?
\end{exercise}

\begin{exercise}[37]
Докажите, что уравнение $X^{3} + Y^{3} + Z^{3} = 2$ имеет бесконечно много решений в целых числах.
\end{exercise}

\begin{Solution}
Рассмотрим тождество
$$
(1+x)^{3} + (1-x)^{3} - 6x^{2} = 2
$$
Тогда положив $x = 6s^{3}$ получим параметризацию
$$
(1+6s^{3})^{3} + (1-6s^{3})^{3} - (6s^{2})^{3} = 2
$$ 
\end{Solution}

\chapter{1972}

\begin{exercise}[33]
Подряд написано 99 девяток. Докажите, что справа к ним можно приписать 100 цифр так, чтобы получившееся число было точным квадратом.
\end{exercise}


\chapter{1973}

\begin{exercise}[10]
Докажите, что $2^{10} + 5^{12}$ - составное число.
\end{exercise}

\begin{Solution}
Начнем с тождества
$$
n^{4} + 4m^{4} = (n^{2} + 2m^{2} + 2nm)(n^{2} + 2m^{2} - 2nm)
$$
А значит $n = 125$, $m = 4$ и
$$
2^{10} + 5^{12} = 125^{4} + 4 \cdot 4^{4} = (125^{2} + 2 \cdot 4^{2} + 2 \cdot 125 \cdot 4)(125^{2} + 2 \cdot 4^{2} - 2 \cdot 125 \cdot 4) = 16657 \cdot 14657
$$
\end{Solution}

\begin{exercise}[37]
Играют двое. Один задумывает десятизначное число, а второй может спрашивать у него о том, какие цифры стоят на конкретном наборе мест в записи. Первый отвечает ему, но без указания того, какие именно цифры на каких именно местах стоят. За какое минимальное число вопросов заведомо можно отгадать число?
\end{exercise}


\chapter{1974}

\begin{exercise}[25]
Существует ли 20-значное число, являющеееся точным квадратом, десятичная запись которого начинается с 11 единиц?
\end{exercise}

\chapter{1975}

\begin{exercise}[36]
Из вершин правильного 35-угольника выбрали восемь. Докажите, что какие-то четыре из них образуют трапецию или прямоугольник.
\end{exercise}


\chapter{1978}

\begin{exercise}[14]
Пятизначное число делится на 41. Докажите, что любое пятизначное число, полученное из него круговой перестановкой цифр, также делится на 41.
\end{exercise}

\begin{exercise}[28]
Клетки доски $100 \times 100$ окрашены в четыре цвета, причем клетки, окрашенные в один цвет, не имеют общих вершин. Докажите, что клетки в углах доски окрашены в разные цвета.
\end{exercise}

\chapter{1979}

\begin{exercise}[22]
Дано $n$ попарно взаимно простых натуральных чисел $a_{1}, a_{2}, \dots, a_{n}$, причем $1 < a_{i} < (2n - 1)^{2}$. Докажите, что среди них есть хотя бы одно простое.
\end{exercise}


\chapter{1982}

\begin{exercise}[23]
Клетки прямоугольника $5 \times 41$ раскрашены в два цвета. Докажите, что можно выбрать три строки и три столбца так, что все девять клеток, находящиеся на их пересечении, будут иметь один цвет.
\end{exercise}

\begin{Solution}
В любом столбце какой-то цвет будет занимать не менее трёх клеток. Нас будет интересовать множество позиций, занимаемых тремя клетками одного цвета. Возможных троек позиций в столбце всего $\binom{5}{3} = 10$.
С другой стороны, найдется не менее 21 столбца с одним преобладающим цветом. Так как возможных троек всего 10, то по принципу Дирихле, cреди этих столбцов найдутся по крайней мере три, в которых клетки этого цвета расположены одинаково.
\end{Solution}

\begin{exercise}[44]
Докажите, что множество всех натуральных чисел, больших единицы, нельзя разбить на два непустых множества так, что если числа $a$ и $b$ лежат в одном множестве, то и число $ab - 1$ лежит в том же множестве.
\end{exercise}

\chapter{1983}

\begin{exercise}[17]
Докажите, что число $2^{58} + 1$ можно представить как произведение трех натуральных чисел, больших 1.
\end{exercise}

\begin{exercise}[31]
Две последовательности чисел $(x_{n})$ и $(y_{n})$ построены так, что $x_{1} = x_{2} = 10$, $y_{1} = y_{2} = -10$ и
$$
x_{n+2}=(x_{n}+1)x_{n+1}+1, \quad y_{n+2}=(y_{n+1} + 1)y_{n}+1
$$
при всех натуральных $n$. Докажите, что любое натуральное число встречается не более чем в одной из этих последовательностей.
\end{exercise}

\begin{exercise}[37]
В правильном 20-угольнике отметили девять вершин. Докажите, что найдется равнобедренный треугольник с вершинами в отмеченных точках.
\end{exercise}

\begin{Solution}
Разделим вершины 20-угольника на 4 множества так, что каждое множество является правильным пятиугольником. В каком-то из пятиугольников, по принципу Дирихле, будет отмечено 3 вершины. Но любые три вершины в пятиугольнике образуют равнобеденный треугольник.
\end{Solution}

\chapter{1984}

\begin{exercise}
Множество $A$ состоит из натуральных чисел, причем среди любых 100 идущих подряд натуральных чисел есть число из $A$. Докажите, что в $A$ найдутся четыре различных числа $a$, $b$, $c$, $d$ такие, что $a + b = c + d$.
\end{exercise}

\chapter{1987}

\begin{exercise}[42]
Разложите число $989 \cdot 1001 \cdot 1007 + 320$ на простые множители.
\end{exercise}

\chapter{1990}

\begin{exercise}[9]
На экране компьютера - число 123. Компьютер каждую минуту прибавляет к числу на экране 102. Программист Федя в любой момент может изменить число на экране, переставив произвольным образом его цифры. Может ли Федя действовать так, чтобы на экране всегда оставалось трехзначное число?
\end{exercise}

\begin{Solution}
Да. Можно зациклить процесс следующим образом
$$
123 \rightarrow 225 \rightarrow 327 \rightarrow 429 \rightarrow 531 \rightarrow 135 \rightarrow 237 \rightarrow 327 \rightarrow \dots 
$$
\end{Solution}

\begin{exercise}[46]
Рассмотрим все возможные наборы чисел из множества $\{1, 2, \dots, N\}$, не содержащие двух соседних чисел. Докажите, что сумма квадратов произведений чисел в этих наборах равна $(N + 1)! - 1$.
\end{exercise}

\chapter{1991}

\begin{exercise}[36]
Можно ли разбить числа $1, 2, 3, \dots, 100$ на три группы так, чтобы в первой группе сумма чисел делилась на 102, во второй - на 203, а в третьей - на 304?
\end{exercise}

\begin{exercise}[42]
Дано 70 различных натуральных чисел, не превосходящих 200. Докажите, что какие-то два из них различаются на четыре, пять или девять.
\end{exercise}

\begin{Solution}

\end{Solution}

\begin{exercise}[47]
В ряд выписано 26 ненулевых цифр. Докажите, что этот ряд можно разбить на несколько частей так, чтобы сумма чисел, образованных цифрами каждой из частей, делилась на 13.
\end{exercise}

\begin{exercise}[53]
С натуральным числом разрешается делать следующие две операции:
\begin{itemize}
\item умножать его на любое натуральное число;
\item вычеркивать в его десятичной записи нули.
\end{itemize}
Докажите, что при помощи этих операций из любого числа можно получить однозначное число.
\end{exercise}

\begin{exercise}[56]
Докажите, что число $512^{3} + 675^{3} + 720^{3}$ - составное.
\end{exercise}

\begin{Solution}
Обозначим $a = 512 = 2^{9}$, $b = 675 = 3^{3}5^{2}$, $c = 720 = 2^{4}3^{2}5$.
\end{Solution}



\chapter{1992}

\begin{exercise}[19]
Клетки квадрата $7 \times 7$ раскрашены в два цвета. Докажите, что найдется по крайней мере 21 прямоугольник с вершинами в центрах клеток одного цвета и со сторонами, параллельными сторонам квадрата.
\end{exercise}

\begin{Solution}
Рассмотрим триплет $(c, p_{1}, p_{2})$, где $c$ это цвет клетки, а $p_{1}$ и $p_{2}$ это две возможные позиции в строке. Всего таких триплетов $2 \cdot \binom{7}{2} = 42$. Теперь посчитаем сколько триплетов покрывает одна строка. Для строки, содержащей $k$ клеток одного цвета, это
$$
\binom{k}{2} + \binom{7 - k}{2}
$$  
Несложно проверить, что минимум этого выражения равен 9 и достигается при $k = 3$ или $k = 4$. Значит все строки покрывают в сумме не менее $9 \cdot 7 = 63$ триплета. Так как возможных триплетов только 42, то повторяющихся триплетов $63 - 42 = 21$.
\end{Solution}

\begin{exercise}
С натуральным числом $A$ разрешается выполнять следующую операцию: разбить его на два натуральных слагаемых, больших 1, и заменить $A$ на их произведение. Докажите, что из любого числа, большего 4, такими операциями можно получить точную степень десяти.
\end{exercise}

\begin{exercise}[49]
На доске написано 128 единиц. За ход можно заменить пару чисел $a$ и $b$ на число $ab+1$. Пусть $A$ - максимальное число, которое может получиться на доске после 127 таких операций. Какова его последняя цифра?
\end{exercise}


\chapter{1993}

\begin{exercise}[20]
На одном из полей доски $7 \times 7$ стоит фишка. Разрешается последовательно ставить на пустые поля новые фишки, но так, чтобы каждое поле, на которое выставляется очередная фишка, имело общую сторону не более чем с одним уже занятым полем. Какое максимальное количество фишек может оказаться на доске после нескольких таких операций?
\end{exercise}

\begin{exercise}[21]
Найдите все натуральные числа, являющиеся степенью двойки, такие, что из них можно получить другую степень двойки, приписав какую-то цифру спереди.
\end{exercise}

\begin{exercise}[27]
Даны два правильных 10-угольника, в каждой вершине каждого из которых написано какое-то натуральное число. Известно, что сумма чисел на каждом 10-угольнике равна 99. Докажите, что на обоих многоугольниках можно отметить несколько подряд стоящих вершин (может быть, одну, но не все) так, что суммы отмеченных чисел будут одинаковы.
\end{exercise}

\begin{exercise}[33]
В квадратной таблице $17 \times 17$ закрашены в черный цвет 80 клеток, остальные - белые. Разрешается закрасить строку или столбец в черный цвет, если большинство клеток на этой линии - черные. Докажите, что при помощи таких операций нельзя сделать всю таблицу черной.
\end{exercise}

\begin{exercise}[45]
Докажите, что квадрат нельзя разрезать на равнобедренные треугольники с углом при вершине, равным $10^{\circ}$.
\end{exercise}

\begin{exercise}[56]
В вершинах правильного $n$-угольника расставлены числа: $n - 1$ нулей и одна единица. Разрешается увеличить на 1 все числа в вершинах любого правильного $k$-угольника, вписанного в данный многоугольник. Можно ли такими операциями сделать все числа равными?
\end{exercise}

\begin{exercise}[64]
По контуру каждой грани выпуклого многогранника ползает муха (таким образом, мух столько же, сколько граней), и все они двигаются, обходя свою грань по часовой стрелке. Известно, что их скорости в любой момент времени не меньше 1 мм/ч. Докажите, что рано или поздно какие-то две мухи столкнутся.
\end{exercise}

\chapter{1994}

\begin{exercise}
В клетках квадратной таблицы $10 \times 10$ расставлены 0 и 1, причём известно, что из любых четырех строчек таблицы какие-то две совпадают. Докажите, что в таблице есть два одинаковых столбца.
\end{exercise}

\begin{exercise}
Есть сто монет достоинством $1, 2, 3, \dots, 100$ пиастров. Среди них не более 20 фальшивых, то есть таких, что их вес в граммах не равен их достоинству. Как при помощи чашечных весов без гирь определить, фальшива ли монета достоинством в 10 пиастров?
\end{exercise}

\chapter{2023}

\chapter{2024}



\end{document} 
 
 
